\section{An actual product-torus family and its full external-row zero tail}
\label{sec:r12-external-row-torus-tail}

\paragraph{Status and attribution.}
\textbf{PROVED THIS ROUND, submitted for exact-source review} is the
complete family and local-uniform zero-tail certificate below.
The construction contributors are \texttt{/root} and the R12 matrix
agent \texttt{/root/r11\_challenge}.  Constructing a primitive by
integration of this family requires the separate, precisely stated
product-torus actual homotopy interface.  No arbitrary proper
pushforward of actual forms or descent from a field extension is
assumed here.

\subsection{The original-field parameter action}

Use the exact ECR construction with the original block $S$ in
rows $2,3$, external row $1$, original tuple $g$, and original
analytic row vector $v=(v_2,v_3)$.  Work on the ordinary overlap
of the two Gaussian pivots of $S$, and write $W$ for its complement
of $|g=0|$.  Let $t$ be the coordinate of an \emph{additional}
original-$K$ torus $\mathbf G_m$, and set
\[
 \tau=v(t),\qquad g^t=D_1(t)g,
 \qquad v^t=t v,\qquad
 V(v^t)g^t=D_1(t)V(v)g.
 \tag{EFT1}
\]
All block coefficients and the original metric are pulled back
from the base.  The common zero set of $g^t$ is $|g=0|$ times
$\mathbf G_m$.  The prescribed ECR formulas therefore define
actual forms on the whole $W\times\mathbf G_m$:
\[
 \begin{split}
 \mathcal T^t={}&q_\alpha(tv;D_1(t)g)-q_\gamma(tv;D_1(t)g)\\
 &-Q_S(D_1(t)V(v)g)+Q_S(D_1(t)g).
 \end{split}
 \tag{EFT2}
\]
The ECR proof applies on this product analytic space, not merely
on each numerical valuation slice.  Consequently
\[
                   \partial_2\mathcal T^t=0
                  \quad\hbox{in }B(W\times\mathbf G_m),
 \qquad (t=1)^*\mathcal T^t=\mathcal T(S,v;g).
 \tag{EFT3}
\]
The first identity includes every prime and double-prime parameter
component.  The second is the original-field analytic section
$t=1$ and uses the actual pullback theorem, not evaluation of
positive-degree forms at a Gauss point.

\begin{theorem}[Full locally uniform ordinary zero tail]
For every $x\in W$ there are an ordinary base neighborhood
$O_x\subset W$ and a finite real constant $T_x\ge0$ such that
\[
            \mathcal T^t=0
       \quad\hbox{on }O_x\times\{\tau>T_x\}
                         \text{ as an actual form}.
 \tag{EFT4}
\]
On every bounded real interval $I\subset\mathbf R$, the family
is locally represented on $\tau\in I$ by a finite product
presentation in base units and the independent parameter
unit $t$.  Every unit occurring in its nonzero expression is of
the form $t^a u$, where $a\in\mathbf Z$ and $u$ is a base unit.
On each cell its smooth coefficients are polynomial in $\tau$,
with coefficients from one finite base smooth presentation.
No unit $1+t u$, no independent change to an analytic sum, and
no arbitrary translation of a native analytic-zero presentation
is introduced.
\end{theorem}

\subsection{Every unpartitioned crossing term vanishes in the tail}

For a fixed Gaussian word let $g_j,v^{(j)}$ be ECR6.  Since each
block factor fixes row $1$, the family gives exactly
\[
       g_j^t=D_1(t)g_j,\qquad
                  (v^{(j)})^t=t v^{(j)}.
 \tag{EFT5}
\]
Thus every external shear coefficient in the complete ECR4
rewrite algorithm is $t a$ for a base analytic function $a$.
Every block coefficient and unit remains a base function.
The following table records a sufficient \emph{whole-witness}
zero domain for every kind of term in that algorithm:
\[
\begin{array}{p{0.44\linewidth}p{0.45\linewidth}}
\text{Full native call}&\text{Sufficient strict inequalities}\\\hline
\text{NSC, external coefficient }ta\text{ and source unit }\lambda
 &|ta|<1,\quad |ta/\lambda|<1;\\
\text{Passive MX, external coefficient }ta
 &|ta|<1;\\
\text{DTC, external coefficient }ta\text{ and block coefficient }c
 &|ta|<1;\\
\text{ST, external coefficient }ta\text{ and block coefficient }c
 &|ta|<1,\quad |tac|<1;\\
\text{SAR, external coefficients }ta,tb,t(a+b)
 &|ta|<1,\quad |tb|<1;\\
\text{Same-target SAN commutation, coefficients }ta,tb
 &|ta|<1,\quad |tb|<1;\\
\text{Signed permutation}&\text{Its prescribed native pair is }0.
\end{array}
 \tag{EFT6}
\]
For NSC, passive MX, DTC, ST and SAR these are exactly the
ordinary zero domains of their full frozen triples, not merely
zero boundaries.  In the SAN row, every individual external
shear companion is zero on its small-coefficient open by TCT2;
their sum is therefore zero, including the complete MT part.
The signed ST insertion ECR5 introduces only $ta$ and $tac$
among the new external coefficients and the unchanged block
coefficient $\pm c$.  Its added inverse pairs cancel literally.

There is thus a finite list $\mathcal B$ of base analytic
functions obtained from the two fixed Gaussian words: include
$a,a/\lambda$ for each NSC, $a$ for passive MX and DTC,
$a,ac$ for ST, and $a,b$ for SAR and same-target commutation.
All denominators here are the already specified base analytic
units on the pivot overlap.  Define
\[
                       A_0=\max(1,|b|:b\in\mathcal B).
 \tag{EFT7}
\]
On the whole ordinary open $|t|A_0<1$, every term in both
crossing triples is zero by (EFT6), for every input tuple:
\[
             q_\alpha(tv;D_1(t)g)
                   =q_\gamma(tv;D_1(t)g)=0.
 \tag{EFT8}
\]
This includes zeros of external or block coefficients.  The
finite function $A_0$ is locally bounded; hence (EFT8) has a
constant tail threshold on a neighborhood of each base point.

\subsection{Full GB norm and cap equality, including section zeros}

Put
\[
             w=v_2g_2+v_3g_3,\qquad
             R_0=\|(g_2,\ldots,g_r)\|,\qquad
             M_1=\max(\|g_1\|,\|g_1+w\|).
 \tag{EFT9}
\]
The norms are taken in the same original metric, so all comparisons
below are independent of a local trivialization.  The two inputs
to $Q_S$ in (EFT2) differ only in row $1$, from $tg_1$ to
$t(g_1+w)$.  Every analytic matrix in the exact GB10/GB16
calculation acts only in rows $2,3$ and fixes row $1$.

Fix a base point in the pivot overlap.  The whole GB call list
can be considered on an ordinary neighborhood with the following
finite description.  If a coefficient is zero, its entire frozen
zero-domain call is deleted as an actual zero form on an ordinary
base neighborhood; no singular summand is deleted separately.
All retained diagonal normalizers and their inverses are base
analytic matrices.  There are then finitely many full norm
vertices $F_\ell g$ and finitely many capped passive norms of
the form
\[
 \max\bigl(\|P_\ell F_\ell g\|,
                       \lambda_\ell\|H_\ell g\|\bigr),
                     \qquad \lambda_\ell>0.
 \tag{EFT10}
\]
Here $P_\ell$ is a fixed subtuple projection, $F_\ell,H_\ell$
are invertible block matrices fixing row $1$, and
$\lambda_\ell=\exp(-\kappa_\ell)$ is the exact positive floor
from that frozen cap, including its fixed positive offset.
The $\kappa_\ell$ depend only on the base block coefficients.
This description includes the $r-1$-entry SL/SAR caps and the
$r-2$-entry CS/NSC/WNR caps separately.  They are not changed
to a new common passive block.

Consider first the ordinary open
\[
                  O_0=\{\|w\|<\|g_1\|\}.
 \tag{EFT11}
\]
The ultrametric identity gives $\|g_1+w\|=\|g_1\|$ there.
All other rows of every corresponding vertex are identical.
Consequently every full norm, every projected subtuple norm,
and every capped norm (EFT10) agrees for the two inputs, for
\emph{all} $t$.  The base coefficient weights and the metric
are also identical.  Thus the two complete $Q_S$ expressions
agree on $O_0\times\mathbf G_m$, including all actual mixed
derivatives and parameter components.

On the other ordinary open $O_1=\{R_0>0\}$ define, using the
finite retained list on the neighborhood under consideration,
\[
 c_0=\min\left(1,\ \|F_\ell^{-1}\|^{-1},
                   \lambda_\ell\|H_\ell^{-1}\|^{-1}
                                    \text{ for all }\ell\right)>0.
 \tag{EFT12}
\]
The matrices act invertibly on the rows other than $1$ as well.
Their remaining-row norm is therefore at least
$\|F_\ell^{-1}\|^{-1}R_0$, or
$\|H_\ell^{-1}\|^{-1}R_0$, respectively.  On the ordinary open
\[
                        |t|M_1<c_0R_0
 \tag{EFT13}
\]
the changing first row is strictly below every full norm floor
and every positive cap floor.  The full norms for the two inputs
are then determined by identical remaining rows.  The floor
branches in (EFT10) are identical; any changing first-row part
of its projected subtuple is strictly below that floor.
Hence the entire capped norms agree too.  This is the reason
one must retain the positive floors: equality of full norms
alone would not control an uncapped passive subtuple.

The two ordinary opens $O_0,O_1$ cover $W$.  Indeed, if $R_0=0$
then $g_1\ne0$, and $w=0$, so the point lies in $O_0$.
On $O_1$ the positive denominator $c_0R_0$ is locally bounded
away from zero after shrinking about the given point, whereas
$M_1$ is locally bounded.  Thus (EFT13) holds for all
$\tau>T_x$ on an ordinary base neighborhood, with a finite
constant $T_x$.  On $O_0$ it is unnecessary.  Combining these
equalities with (EFT8) proves (EFT4).

All arguments used equalities on ordinary product opens.
They retain partial zeros of the section tuple, including
$g_2=g_3=0$, zeros of arbitrary other passive rows, and zeros
or cancellations in $g_1+w$.  No assertion is obtained by
testing only a strict Gauss chamber or a dense unit locus.

\subsection{The exact finite parameter presentation}

We detail why this family has the required product presentation.
On any local trivialization and any bounded parameter slab
$\tau\in I$, with $I\subset\mathbf R$ a bounded interval,
all nonzero base analytic quantities used
below can be included in one finite original-$K$ unit chart.
At a zero quantity, either its entire call is zero on the
ordinary coefficient-zero neighborhood (shrunk using the finite
upper bound $\sup_{\tau\in I}|t|$ for an external coefficient),
or a sufficiently
small original-$K$ constant replaces a zero norm entry on an
ordinary neighborhood without changing its dominating norm
or positive floor.  Both a positive lower bound and a finite
upper bound for $|t|$ hold on this bounded slab, so the latter replacement can
be made uniformly on that slab.  This makes no assertion
about replacing an individual zero section uniformly out
to $\tau=+\infty$; there the whole family is already zero
locally by (EFT4).

For every retained term the exact parameter dependence is:
\[
\begin{array}{p{0.39\linewidth}p{0.50\linewidth}}
\text{Part of the native construction}&\text{Parameter dependence}\\\hline
\text{Block coefficients, pivots, determinant, original metric}
 &\text{Pulled back from the base.}\\
\text{External coefficients and their SAR sums}
 &ta,\ tb,\ t(a+b)\text{ with base }a,b.\\
\text{External ST product coefficient}
 &tac\text{ with base }a,c.\\
\text{Input and all block-prefix tuples}
 &\text{First row }t h_1;\text{ all other rows }h_j.\\
\text{An external SHT midpoint normalizer}
 &D_1(ta)^{-1}(t h_1,h_2,\ldots)
       =(h_1/a,h_2,\ldots);\text{ no mixed }t\text{ sum}.\\
\text{Block midpoint normalizers and signed permutations}
 &\text{Base matrices; first-row exponent unchanged.}\\
\text{Weights and all native caps}
 &\text{Finite affine valuations }a\tau+v(u)
   \text{ and their prescribed finite PL minima/maxima.}
\end{array}
 \tag{EFT14}
\]
In the external SHT step its prescribed $E_{1j}(1)$ isometry
acts only after the common $t$ factor has been removed from
row $1$, so its first-row sum is a sum of base analytic
functions.  The final $D_1(ta)$ restores the same $t$ factor.
The signed ST insertion and all same-target collection steps
retain this property.  Every norm vertex can therefore be
written with row exponents $0$ or $1$, never as a new unit
$1+t a$.  Inverses of monomial units only change their integer
exponent.  The two GB caps in the last difference of (EFT2)
remain exactly their old caps and have base-only floor functions.

The functions in the last row of (EFT14) have one finite rational-
direction polyhedral presentation with all real intercepts
retained.  CS, SHT and MT are finite polynomial expressions
in their potentials, unit weights, original metric and their
derivatives.  The only nonpolynomial smooth coefficients here
are the original base metric coefficients and the base-only GB
partition.  Thus, after subdivision, every parameter coefficient
is polynomial in $\tau$ with smooth base coefficients.  The
full differential includes both parameter colours and their
mixed carrier terms; (EFT3) is an identity for this entire
actual presentation.

For any base affinoid and any base unit refinement, adjoining
the independent torus coordinate $t$ gives the product
tropicalization with a multiplicity-one interval away from
the annulus boundary.  Every monomial unit $t^a u$ is an
integral affine function of these product coordinates.
This is a consequence of the original-field product torus,
not an assertion about an arbitrary many-to-one pushforward.
It supplies exactly the product-cycle data needed by an
actual product-parameter homotopy proof.  Possible further
refinements involving other analytic units are not discarded:
the input closedness already holds as an actual form on
the product, and any proposed output must be verified on
every base refinement by the actual zero criterion.

\paragraph{What is and is not concluded here.}
The family (EFT2) is a full actual closed triple, restricts to
the desired ECR7 triple at $t=1$, has a locally uniform
ordinary zero tail, and satisfies the explicit finite product
presentation conditions.  These facts are now proved, not
assumed as a hypothetical contraction of a free presentation.
If the separate product-presented Schweitzer homotopy is
proved for precisely these hypotheses, it can be applied
without a field extension or an unspecified actual-form
pushforward.  That additional theorem, not just the limit
of the pointwise coefficients, is what must supply an actual
$\partial_3$ primitive.  No primitive or general matrix
coherence theorem is asserted by the tail estimate alone.

\paragraph{Exact inputs.}
ECR: \path{33afe064e462e14025981051752eed387264c52f199076b19cc45c4440d98552}.
GB: \path{f4f87d6aabfc325f7d2f01d629a6f4a9fa80227cb9600cf592ea380b0b5104d9}.
GIC: \path{2ea4db21676b04cc4f2bbebfe168fec80c405f228e546f16550b9445d707c056}.
All exact ST/DTC/MX/NSC/SAR/SAN/SHT/MT sources are those listed
in ECR, SAN and TCT.  Mihatsch's actual pullback and all-refinement
zero criterion are Theorem 4.5 and Lemma 4.4 of
\texttt{arXiv:2112.10018v4}; no stronger pushforward statement
is attributed to them.
